%%%
%%% S12: Region-calibrated personas — preregistered null result
%%% CogSci 2027 submission (anonymized)
%%%

\documentclass[10pt,letterpaper]{article}

\usepackage{cogsci}
\usepackage{graphicx}
\usepackage{booktabs}
\usepackage{amsmath}
\usepackage[round]{natbib}

% \cogscifinalcopy %%% Uncomment for camera-ready

\title{Region-Calibrated Personas Do Not Significantly Improve Fit to
CPT-Imputed Country-Level Choice Distributions:
A Preregistered Null Result}

\author{{\large \bf Anonymous} \\
\texttt{anonymous} \\
\hspace{1cm}Anonymized for double-blind review}

\begin{document}

\maketitle

\begin{abstract}
Language models are increasingly used as stand-ins for human participants
in behavioral research, yet they are systematically \textit{more} rational
than the people they simulate. One proposed fix infuses calibrated
human-like irrationality at inference time; the open question is whether
that calibration should be region-specific, given real cross-cultural
variation in risk preferences. We preregistered exactly that test: do personas
calibrated to \textit{region-specific} measured preference distributions
reproduce country-level choice heterogeneity better than a generic-human
calibration? Running Qwen3-1.7B on CPU, we simulated individual respondents
for ten countries spanning the empirical loss-aversion distribution,
conditioning only on distribution-sampled behavioral profiles---no country
label, ethnicity term, or trait inventory ever appears in a prompt---and
compared simulated choice distributions against prospect-theory-imputed
choice distributions (deterministic predictions from reconstructed
within-country parameter distributions, not observed human choices) on six
held-out willingness tasks. The preregistered primary
test (paired permutation over countries, $\alpha=0.05$) \textbf{rejects the
hypothesis}: region-conditioned simulations average Wasserstein-1 distance
0.304 versus 0.285 for generic calibration (difference $-0.019$; one-sided
$p=0.631$), while the uncalibrated arm (0.294) matches about as well as
either calibrated arm. This null concerns one small
model under fixed protocols---not language models in
general---and we discuss what it does and does not rule out for
pluralistic alignment.

\textbf{Keywords:}
cultural calibration; language models; behavioral simulation; null results;
preregistration; prospect theory
\end{abstract}

\section{Introduction}

Language models are being recruited as simulated human subjects: they answer
surveys, play economic games, and stand in for participants in market
research \citep{argyle2023,aher2023,park2023}. The practice rests on an
untested premise: that a model's responses reproduce the decision patterns
of the humans it replaces. The premise fails, and it fails in a specific,
replicable direction. Across cognitive-reflection batteries, lottery choices, and
repeated games, models deviate from humans toward \textit{greater}
rationality: they are too accurate \citep{aher2023}, too close to
expected-value maximization \citep{liu2025}, and too consistent to be human
\citep{macmillan-scott2024}. Instruction tuning can erase human-like biases
that smaller base models display \citep{hagendorff2023}. Call this the
\textit{rationality gap}.

If the gap is the disease, the natural treatment is deliberate: install
calibrated irrationality by prompting or steering the model until its
measured biases match human effect sizes. Early attempts show persona
conditioning can elicit prospect-theory signatures \citep{dubedy2026} and
motivated reasoning in models---yet none calibrates the installed
irrationality against measured human distributions. That generic-calibration
problem is taken up in a companion paper. Our starting observation: even a perfectly calibrated \textit{generic}
human is the wrong target, because human irrationality is not a single
profile. Risk
preferences vary substantially across countries \citep{rieger2015}, loss
aversion varies with culture \citep{wang2017}, and the WEIRD-psychology
literature shows that Western convenience samples are outliers on exactly
the traits models default to \citep{henrich2010}. A simulator calibrated to
a single human profile flattens this heterogeneity.

The question we ask is narrower, and falsifiable: does calibrating
irrationality levers \textit{region-specifically}---to measured cross-country
preference distributions---reproduce cross-country heterogeneity in
economic choice better than a generic-human calibration? The ground truth is
the INTRA survey (International Test on Risk Attitudes), which elicited risk
preferences from roughly 7,000 respondents in 53 countries
\citep{rieger2015}, with per-country cumulative-prospect-theory parameters
estimated in follow-up work \citep{rieger2017}. To be precise about the sense in which
this is ``ground truth'': the choice distributions we fit are not observed
human choices but deterministic CPT predictions computed from
\textit{reconstructed} within-country parameter distributions
(within-country SDs recovered from published SEs under an assumed $n=150$,
winsorized). What we measure is therefore how well qualitative personas
track a parametric imputation of human heterogeneity---twice removed from
actual choices---not direct human-likeness. The full
design---countries, levers, tasks, metrics, tests, and decision
rules---was preregistered before any model ran, and we report the
preregistered test exactly as specified, null outcome included.

Three things come out of this work. First, a method for
distribution-grounded regional calibration that satisfies a hard
anti-essentialism constraint:
levers are calibrated to measured \textit{distributions} with
within-country variance, prompts contain individual behavioral profiles
rather than country labels, and culture is treated as situated and
contextual \citep{oyserman2008} rather than as fixed national character.
Second, a preregistered empirical test---ten countries spanning the
measured loss-aversion distribution, region-conditioned versus generic-human
versus uncalibrated simulation on held-out willingness tasks, fit against
prospect-theory-imputed choice distributions---that returns a null. Third, a constraint on the pluralistic-alignment agenda. Region-specific
calibration, as implemented here, does not buy country-level fit. Along the
way, a cautionary observation on response format: the Yes/No willingness
format flips the usual rationality-gap pattern relative to the
prospect-theory prediction. We report that flip as format
sensitivity---not as a challenge to the human-data finding---and it limits
what ``human-like'' can mean as a calibration target.

\section{Related work}

\subsection{The rationality gap between LLMs and humans}

Excess rationality is the most reliable fidelity failure in LLM-as-human
simulation. \citet{aher2023} named it \textit{hyper-accuracy distortion}:
aligned models are too factually correct to reproduce human response
distributions. The direct decision-making statement comes from
\citet{liu2025}: against the 13,000-choice human dataset choices13k,
GPT-4o, Claude 3 Opus, and Llama-3 models deviate from humans toward
expected-value theory, both when simulating and when predicting choices.
Alignment itself is partly responsible---\citet{hagendorff2023} show
mid-scale GPT-3 gave the intuitive wrong answer on cognitive-reflection
items about 80\% of the time (humans about 55\%), while instruction-tuned
ChatGPT scored 59--96\% correct. The settled picture: models are
\textit{differently} biased rather than uniformly biased
\citep{pagliaro2026,suri2024}, and bias is a prompt--model interaction
property, not a stable trait.

\subsection{Attempts to make models irrational}

Prompting can nudge models toward human-like bias, but without calibration.
LLM ``irrationality'' is inconsistent across repeats and often non-human in
form \citep{macmillan-scott2024}. Persona assignment does elicit motivated
reasoning \citep{dash2025} and prospect-theory signatures---\citet{dubedy2026}
reports GPT-4.1 with income personas reproducing the fourfold pattern
unprompted. The closest concurrent work treats LLMs as calibrated
measurement instruments for behavioral parameters: \citet{yee2026} induce
large, stable shifts in loss aversion, herding, extrapolation, and anchoring
toward human benchmark magnitudes across 24,000 agent--scenario pairs via
profile-based calibration. Their calibration targets generic human
benchmarks, however, not region-specific distributions, and they do not test
whether calibration reproduces cross-country heterogeneity. We leave the
generic-calibration problem to a companion paper and ask instead whether
calibration should be regional.

\subsection{Cross-cultural variation in economic preferences}

Regional variation in irrationality is among the best-established facts in
behavioral economics. Across 15 small-scale societies, \citet{henrich2001}
document dramatic cross-society variation in ultimatum, dictator, and
public-goods play; \citet{henrich2010} argue Western samples are frequent
outliers. On risk, \citet{rieger2015} surveyed 53 countries and found risk
aversion in gains and risk seeking in losses everywhere---but at degrees
that differ significantly across countries. Per-country
cumulative-prospect-theory parameters estimated by \citet{wang2017} show
loss aversion varying systematically with culture. The largest validation is
the Global Preference Survey: experimentally validated preference measures
for 80,000 individuals in 76 countries, with substantial cross-country
heterogeneity---and larger within-country heterogeneity \citep{falk2018}.
That literature gives us our design rule directly: calibrate to
distributions, never to country means \citep{falk2018}.

\subsection{Simulating regional humans, and the stereotype trap}

Country labels do shift model responses toward the target country
\citep{durmus2023,tao2024}---but they invite exactly the failure this paper
is designed to avoid: \citet{durmus2023} show country-perspective prompting
can \textit{amplify} cultural stereotypes, and \citet{cheng2023a} and
\citet{cheng2023b} document stereotyped portrayals and caricature in
persona-conditioned generations. The constructive alternatives are all
distribution- or retrieval-grounded. \citet{abels2026} build ``calibrated
value personas'' from \textit{individual-level} survey responses;
\citet{seo2025} retrieve per-individual value summaries instead of relying
on model priors. We combine both ideas: individual-level profiles sampled
from measured country distributions, and---critically---no country label
anywhere in the prompt, so the model cannot fall back on training-data
national associations. Our theoretical guardrail is \citet{oyserman2008}:
cultural influence is situated and moment-to-moment rather than a static
national trait, so each prompt describes one respondent's measured survey
answers, never a national character.

\section{Method}

We wrote the full preregistration---countries, levers, tasks, metrics,
tests, decision rules, anti-essentialism rules, and two amendments---before
any model ran. What follows summarizes the design \textit{as run},
amendments included.

\subsection{Ground truth: the INTRA survey}

The ideal ground truth would be the Global Preference Survey
\citep{falk2018}, but its microdata require a data-request form and forbid
redistribution, so we excluded it per the preregistration and fell back to
the pre-specified alternative: the INTRA survey \citep{rieger2015}, which
elicited willingness-to-pay for gain lotteries and willingness-to-accept for
loss lotteries from about 7,000 economics students in 53 countries, plus
the per-country cumulative-prospect-theory parameters estimated from it
\citep{rieger2017}---loss aversion $\lambda$, curvature $\phi$,
probability-weighting $\delta$. Numerical values are the published
per-country estimates as tabulated in \citet{xie2018} (open access).

The INTRA publications report per-country standard errors rather than
individual-level data, so we reconstruct within-country distributions
parametrically: each individual's $(\lambda_i,\phi_i,\delta_i)$ is drawn
from a truncated normal centered on the published country estimates, with
$\mathrm{SD}=\mathrm{SE}\times\sqrt{n}$, $n=150$ (median INTRA per-country
sample size), winsorized to published ranges and $\lambda$ truncated at
0.02 (Iceland's published $\lambda=-0.112$ becomes small-positive draws;
disclosed). Absolute fit levels depend on this reconstruction, but the
region-versus-generic comparison does not---both arms draw from the same
reconstructed family---and sensitivity analyses scaling all within-SDs by
0.5 and 1.5 leave the ordering unchanged.

\subsection{Country selection}

We selected ten countries stratified across the empirical loss-aversion
distribution ($\lambda$ is the prospect-theory parameter shown to vary most
with culture, \citet{wang2017}): Iceland ($-0.11$), Greece ($0.33$), Brazil
($0.60$), Hungary ($0.93$), Pakistan ($1.20$), Australia ($1.65$), Germany
($1.90$), Japan ($1.93$), Denmark ($2.21$), Italy ($3.35$). Nothing but
these published numbers drove the selection.

\subsection{Elicitation and calibration levers (Amendment 2)}

Pilot diagnostics forced an elicitation redesign before any main-run
generation (preregistered Amendment 2). The chat template produced
degenerate label bias---near-deterministic single-label outputs regardless
of profile---which raw completion removes. Binary A/B choices proved
content-insensitive to numerical certainty-equivalent profiles: first-token
choice probabilities came out near-identical for Italy-like versus
Hungary-like profiles. The final elicitation is therefore \textbf{raw
completion} (no chat template) with \textbf{qualitative 3-level personas}
and \textbf{Yes/No willingness questions}: e.g.\ ``Would you pay \$45 for a
50/50 chance to win \$100? Answer:'' with the first Yes/No token parsed.

Each simulated individual draws $(\lambda_i,\phi_i,\delta_i)$ either from
their country's reconstructed distribution (\textbf{region} arm) or from
the pooled multi-country mixture (\textbf{generic} arm); the draw is then
mapped to a qualitative persona via absolute tertile cutoffs from the
pooled reconstructed distribution. That mapping yields three risk
sentences (from curvature $\phi$, used on gain tasks) and three loss
sentences (from loss aversion $\lambda$, used on loss tasks), e.g.\ ``You
are very cautious with money. You strongly prefer guaranteed amounts over
any gamble.'' The \textbf{uncalibrated} arm sees the question with no
persona. No prompt in any arm contains a country name, an ethnicity term,
or a trait inventory---a scripted audit verifies this over all prompt
templates and all 1,392 generations (0 hits).

\subsection{Decision tasks}

We held out six Yes/No willingness tasks with prices fixed across
individuals, spanning the human certainty-equivalent range: pay \$25, \$45,
or \$65 for a 50/50 chance to win \$100 (gain); accept a 50/50 chance of
losing \$100 if paid \$40, \$90, or \$140 (loss). Each simulated
individual's ``real'' counterpart is the deterministic choice implied by the
$(\lambda_i,\phi_i,\delta_i)$ drawn for that individual under cumulative
prospect theory \citep{kahneman1979,tversky1992}: yes on a gain task iff
$\mathrm{price} \le 100\cdot0.5^{\phi_i}$; yes on a loss task iff
$\mathrm{price} \ge 100\cdot(0.5\lambda_i)^{\phi_i}$. We use the
Tversky--Kahneman one-parameter weighting form because the exact form in
\citet{rieger2017} is not public (disclosed assumption).

\subsection{Model and scale}

We ran Qwen3-1.7B (Q4\_K\_M GGUF) with CPU-only inference via llama.cpp
(temperature 0.7, top\_p 0.8, top\_k 20, max\_tokens 24, n\_ctx 1024---reduced
from the preregistered 64 and 2048 under Amendment~2's elicitation redesign,
which replaced the chat template and its \texttt{/no\_think} suffix with raw
completion). The preregistered scale after Amendment (N\_INDIV $16\to8$,
to bound CPU time) is 8 individuals $\times$ 10 countries $\times$ 2
conditioned arms $\times$ 6 tasks $=$ 960 main generations, plus 48
uncalibrated baseline generations and 384 robustness generations (prompt
paraphrase, seed variant, temperature 0.3, and a loss-aversion-only
ablation, on Germany and Italy): 1,392 generations in total. Seeds are
deterministic (master seed 20261008); A/B option order does not apply in
the Yes/No format.

\subsection{Metrics and preregistered tests}

Our \textbf{primary} metric is the per-country Wasserstein-1 distance
between the simulated per-individual risky-choice distribution and the
prospect-theory-imputed distribution, averaged over the ten countries; the
preregistered test is a one-sided paired permutation test (10 country
pairs, exact $2^{10}$) of $W_1(\mathrm{generic}) - W_1(\mathrm{region}) > 0$
at $\alpha=0.05$. Note that with 10 pairs the exact permutation
distribution has granularity $1/1024$, so the one-sided $p$ is coarse by
construction. \textbf{H1}---region-conditioned calibration yields smaller
distributional distance than generic calibration---is supported iff this
test rejects in the predicted direction; otherwise it stands rejected.
Secondary metrics, each reported with 95\% bootstrap confidence intervals,
are (a) within-country dispersion preservation; (b) between-country
heterogeneity, as the Spearman rank correlation of country means; and (c)
the EV-alignment gap. Every metric is gated on the output-completeness
check (all 1,392 expected generations present).

After the main run we found and fixed two analysis bugs (the script had
never been exercised end-to-end): a country filter that emptied the pooled
generic arm, and a replicate-grouping error on the generic arm's 8-persona
$\times$ 10-replicate structure. Independent review then surfaced a deeper
problem no analysis fix could cure. The generic arm had reused the same 8
personas identically across all 10 countries, so its per-country $W_1$
estimates shared one simulated distribution by construction---correlating
the paired differences and leaving the generic arm with zero
between-country variation by design. Since this is a data-structure defect,
we re-ran the generic arm with 8 fresh independent draws from the pooled
mixture \textit{per country} (80 individuals total;
\texttt{code/run\_generic\_v2.py})---the reading of the preregistration's
``drawn from the POOLED 27-country distribution'' that gives each country
an independent generic estimate at the same $n$ as the region arm.
Everything else (elicitation, tasks, model, sampling parameters) is
identical; only the persona-draw keys and generation seeds are per-country.
All generic-arm results below use this corrected data, while the original
shared-persona records are retained unmodified as an audit trail. The
preregistered hypotheses, metrics, and decision rule are unchanged, and a
one-command reproduction package (code, data, verifier;
\texttt{bash reproduce.sh}) accompanies the paper so every number below
can be re-verified independently.

\section{Results}

The corrected dataset contains all 1,392 expected generations (48/48
baseline, 480/480 region + 480/480 corrected generic, 384/384 robustness),
with parse rates (first Yes/No token found) of 93.8\% baseline, 88.3\%
region, and 91.2\% corrected generic. The no-label audit comes back clean:
0 hits across prompt templates, all prompts, and all 1,392 generations.

\subsection{Primary test: H1 rejected}

\begin{table}[t]
\centering
\caption{Per-country Wasserstein-1 distance of simulated to
prospect-theory-imputed choice distributions, by arm. Lower is better.
$n_R$ / $n_G$ = complete region / generic individuals after parse failures
(generic arm: corrected per-country independent draws).}
\label{tab:w1}
\begin{tabular}{lccccc}
\toprule
Country & Region & Generic & Uncalibrated & $n_R$ & $n_G$ \\
\midrule
Iceland & 0.292 & 0.181 & 0.442 & 4 & 3 \\
Greece & 0.312 & 0.451 & 0.421 & 4 & 6 \\
Brazil & 0.333 & 0.222 & 0.400 & 4 & 3 \\
Hungary & 0.200 & 0.167 & 0.400 & 5 & 6 \\
Pakistan & 0.132 & 0.354 & 0.254 & 3 & 4 \\
Australia & 0.346 & 0.229 & 0.254 & 5 & 6 \\
Germany & 0.174 & 0.438 & 0.138 & 3 & 2 \\
Japan & 0.667 & 0.357 & 0.200 & 1 & 7 \\
Denmark & 0.250 & 0.200 & 0.200 & 2 & 5 \\
Italy & 0.333 & 0.250 & 0.233 & 3 & 6 \\
\midrule
Mean & 0.304 & 0.285 & 0.294 & & \\
\bottomrule
\end{tabular}
\end{table}

The three arms land almost on top of each other: mean $W_1$ 0.304 for
region, 0.285 for generic, 0.294 for uncalibrated. The paired difference
(generic $-$ region) is $-0.019$, with region favored in only 3 of 10
countries; the one-sided exact permutation test returns $p=0.631$, far from
$\alpha=0.05$. \textbf{H1 is rejected.} We report this as a null, plainly:
no detectable region advantage, and the uncalibrated arm matches about as
well as either calibrated arm---on this metric, calibration adds nothing
over no persona at all. An exploratory sensitivity check that drops the
three thinnest region cells (Japan, Denmark, Pakistan) agrees
(diff $-0.008$, $p=0.539$ on 7 countries).

\subsection{Secondary results}

\textbf{Robustness.} Across Germany and Italy, none of the four elicitation
variants gives the region variant a consistent edge over the corrected
generic: paraphrase $+0.014$, seed $-0.014$, temperature 0.3 $+0.042$,
loss-aversion-only ablation $+0.205$ (mean $W_1$ differences,
generic $-$ region-variant). Only the ablation shows a sizable gap---on two
countries, so suggestive rather than confirmatory.

\textbf{EV-alignment.} Mean share of ``yes'' answers comes out as 0.62 for
region, 0.48 for generic, 0.27 for uncalibrated, 0.50 for
prospect-theory-imputed, and 0.67 for EV-optimal. The region arm answers at
roughly the EV-optimal rate, the corrected generic arm near the
human-predicted rate, and the uncalibrated arm is strongly yes-averse. That
runs counter to the preregistered baseline expectation---the classic
rationality gap, with the uncalibrated model more EV-rational than humans.
We do \textit{not} read this as a non-replication of \citet{liu2025}. Their
finding rests on actual human choices (choices13k) on lottery tasks; ours
compares a raw-completion Yes/No format against a synthetic CPT baseline.
The flip is a cautionary observation about format sensitivity, not evidence
against the published gap.

\textbf{Dispersion and heterogeneity.} Within-country dispersion
preservation carries no information here: several country cells have
near-degenerate simulated distributions, which blows up the log-ratio
(region 8.2, generic 12.0, with enormous bootstrap CIs). Between-country
heterogeneity is null in both arms---Spearman $\rho=0.09$ ($p=0.81$) for
region, $\rho=0.18$ ($p=0.61$) for generic---with no detectable tracking of
real between-country differences. (Thanks to the corrected per-country
generic draws this comparison is now fair; in the original shared-persona
implementation it was degenerate by construction.)

\section{Discussion}

The preregistered bet fails its own test: region-calibrated personas do
not significantly improve fit to country-level choice distributions, and
the uncalibrated baseline matches about as well as either calibrated arm.
The null has three layers, and each teaches something different.

First, the null itself is clean. Under the corrected generic arm there is
no directional hint of a region advantage: paired difference $-0.019$ (3 of
10 countries favor region), $p=0.631$, and the exploratory thin-cell
sensitivity agrees. An earlier analysis had shown a suggestive $+0.072$
($p=0.090$), but that signal was partly an artifact of the shared-persona
structure---with independent per-country generic draws, it vanishes. A
null, not a near-miss, and we claim nothing more.

Second---and more damaging---is the \textit{uncalibrated} half of the
null. Genuine calibration signal should separate the region arm from the
no-persona baseline; instead the two tie (0.304 vs.\ 0.294). Whatever the
personas contribute gets swamped by the model's decisive answering style:
simulated per-individual distributions come out extreme-heavy (many 0.0 or
1.0 values), so the model delivers verdicts rather than graded
propensities. On this metric, in this format, the persona is nearly inert.

Third, a caution about the rationality gap. The literature's reliable
finding---models more EV-rational than humans \citep{liu2025}---is measured
against actual human choices. Our comparison runs against a synthetic CPT
baseline in a raw-completion Yes/No format, so the flipped pattern
(uncalibrated yes-averse at 0.27 vs.\ predicted 0.50) is an
apples-to-oranges contrast. We report it as evidence that response format
shapes what ``human-like'' even means---a complication for any calibration
agenda built on the gap---rather than as a challenge to the published
finding.

For pluralistic alignment, what we have is a boundary condition, not a
verdict: region-specific calibration \textit{as implemented here}---small
model, qualitative personas, willingness format---does not buy
country-level fit. Concretely, Amendment~2 collapsed continuous
$(\lambda,\phi)$ draws into three fixed qualitative sentences via pooled
tertiles, so the region and generic arms differ only in mixture weights
over identical sentences. The null is therefore joint over the hypothesis
\textit{and} this coarse transmission channel---a finer-grained persona
could in principle carry distributional signal that three sentences
cannot. This does not show that no regional calibration could work. It
shows that the cheapest version does not, and that claims about
``culturally calibrated'' simulators need preregistered tests with this
structure before they are believed.

\paragraph{Concurrent work disclosure.}
Independently, \citet{yee2026} develop profile-based calibration of
behavioral parameters (loss aversion, herding, extrapolation, anchoring)
in LLMs against generic human benchmarks, documenting the same baseline
rationality bias this literature reports; their work appeared on arXiv in
February 2026, before this paper's experiments were run (October 2026).
Where we differ is in testing \textit{region-specific} calibration to
measured cross-country distributions---questions \citet{yee2026} do not
address. A second concurrent audit reaches a compatible verdict at much
larger scale. \citet{chen2026} construct one-to-one ``digital twin''
personas for 67,201 Global Preference Survey respondents across 74
countries and find LLM-simulated responses explaining only about 1.3\% of
individual-level and 4.2\% of country-level variance in human economic
preferences; in ablations, removing the country label from the persona
barely changes alignment, while removing demographic cues degrades it
substantially. That models condition on explicit persona content rather
than on latent country associations is consistent both with our no-label
design and with our result that region-calibrated behavioral profiles add
little over an uncalibrated baseline.

\section{Limitations}

Several limits deserve plain statement. The most serious is cell
thinning: parse failures (11.7\% on the corrected main run) left region
cells with 1--5 complete individuals per country (Japan: $n=1$), and a
$W_1$ distance computed from a single draw is not a distribution
estimate---so the primary test's inputs are degraded at exactly the
noisiest cells, which the paired test accounts for only partially.
Dropping unparseable generations was a forced, non-preregistered choice,
and the missingness is non-random on top of that: failures concentrate on
gain tasks (up to 29\% on the \$65 gain task vs.\ near-zero on loss tasks)
and on the middle risk persona (``somewhat cautious'' parses at 49--65\%
vs.\ 86--100\% for the extreme levels). Cell thinning is therefore not
ignorable---though the exploratory thin-cell sensitivity (dropping Japan,
Denmark, Pakistan) agrees with the main result. A second caveat concerns
the generic-arm correction: while it implements the preregistration's
per-individual pooled draws, it required new data collection after the
original results were seen. It is a reviewer-mandated repair rather than a
preregistered step, and a fully preregistered replication would re-run both
arms. Third, Qwen3-1.7B is a
small model on CPU, and larger models may show different calibration
dynamics. Fourth, the Yes/No willingness format (Amendment 2) elicits more
coarsely than the preregistered lottery-choice format; format effects are
visible in the flipped rationality-gap pattern. Fifth, the ``real'' choice
distributions are deterministic prospect-theory predictions from
reconstructed---not observed---within-country parameter distributions,
twice removed from actual human choices. The within-SD sensitivity check
moves the paired difference only from $-0.018$ to $-0.032$ (ordering
unchanged), but absolute fit levels remain fragile to the reconstruction.
The INTRA samples, sixth, are economics students rather than
representative populations, and within-country distributions are
reconstructed from published standard errors rather than observed
microdata. Finally, the exact probability-weighting form in
\citet{rieger2017} is not public; the Tversky--Kahneman form used here is
a disclosed assumption.

\section{Conclusion}

We preregistered a clean test of a natural hypothesis---that calibrating
LLM personas to region-specific measured preference distributions
reproduces country-level choice heterogeneity better than a generic-human
calibration. Run on Qwen3-1.7B, the test returned a null: with the
corrected per-country generic arm, the paired difference is $-0.019$
($p=0.631$)---no detectable region advantage, and calibration adds nothing
over an uncalibrated baseline (0.304 vs.\ 0.285 vs.\ 0.294). The null's
value is discipline. It tells the pluralistic-alignment agenda what the
cheapest version of regional calibration does \textit{not} buy, under
exactly the conditions a preregistration makes checkable. If culturally
calibrated simulators are to be more than a slogan, they will need to pass
tests built with this structure.

\section*{Declaration of generative AI use}

The author used the AI tool Muse (Muse Spark 1.3, Meta) for drafting,
grammatical editing, code validation, and literature search. The contents
are reviewed, edited, and attested for accuracy by the author.

\begin{thebibliography}{99}

\bibitem[Abels et al.(2026)]{abels2026}
Abels, T., et al. (2026). Improving cross-cultural survey simulation with
calibrated value personas. \textit{arXiv:2605.16193}.

\bibitem[Aher et al.(2023)]{aher2023}
Aher, G., Arriaga, R.~I., \& Kalai, A.~T. (2023). Using large language models
to simulate multiple humans. In \textit{Proc.\ ICML 2023}.

\bibitem[Argyle et al.(2023)]{argyle2023}
Argyle, L.~P., et al. (2023). Out of one, many: Using language models to
simulate multiple humans. In \textit{Political Analysis, 31}(3), 337--351.

\bibitem[Chen et al.(2026)]{chen2026}
Chen, F., et al. (2026). Simulated responses from large language models are
weak proxies for human economic preferences across 74 countries.
\textit{Nature Human Behaviour}.
https://doi.org/10.1038/s41562-026-02611-4

\bibitem[Cheng et al.(2023a)]{cheng2023a}
Cheng, M., Durmus, E., \& Jurafsky, D. (2023). Marked personas: Using
natural language prompts to measure stereotypes in language models.
In \textit{Proc.\ ACL 2023}.

\bibitem[Cheng et al.(2023b)]{cheng2023b}
Cheng, M., Piccardi, T., \& Yang, D. (2023). CoMPosT: Characterizing and
evaluating caricature in LLM simulations. In \textit{Proc.\ EMNLP 2023}.

\bibitem[Dash et al.(2025)]{dash2025}
Dash, S., Reymond, A., Spiro, E.~S., \& Caliskan, A. (2025).
Persona-assigned large language models exhibit human-like motivated
reasoning. \textit{arXiv:2506.20020}.

\bibitem[Dubedy(2026)]{dubedy2026}
Dubedy, S. (2026). Persona-conditioned risk behavior in LLMs.
\textit{arXiv:2603.15831}.

\bibitem[Durmus et al.(2023)]{durmus2023}
Durmus, E., et al. (2023). Towards measuring the representation of subjective
global opinions in language models. In \textit{Proc.\ COLM 2024}.

\bibitem[Falk et al.(2018)]{falk2018}
Falk, A., Becker, A., Dohmen, T., Enke, B., Huffman, D., \& Sunde, U. (2018).
Global evidence on economic preferences. \textit{Quarterly Journal of
Economics, 133}(4), 1645--1692.

\bibitem[Hagendorff et al.(2023)]{hagendorff2023}
Hagendorff, T., Fabi, S., \& Kosinski, M. (2023). Human-like intuitive
behavior and reasoning biases emerged in large language models but
disappeared in ChatGPT. \textit{Nature Computational Science, 3}(10),
833--838.

\bibitem[Henrich et al.(2001)]{henrich2001}
Henrich, J., et al. (2001). In search of homo economicus: Behavioral
experiments in 15 small-scale societies. \textit{American Economic Review,
91}(2), 73--78.

\bibitem[Henrich et al.(2010)]{henrich2010}
Henrich, J., Heine, S.~J., \& Norenzayan, A. (2010). The weirdest people in
the world? \textit{Behavioral and Brain Sciences, 33}(2--3), 61--83.

\bibitem[Kahneman \& Tversky(1979)]{kahneman1979}
Kahneman, D., \& Tversky, A. (1979). Prospect theory: An analysis of decision
under risk. \textit{Econometrica, 47}(2), 263--291.

\bibitem[Liu et al.(2025)]{liu2025}
Liu, A., Geng, M., Peterson, J.~C., Sucholutsky, I., \& Griffiths, T.~L.
(2025). Large language models assume people are more rational than we really
are. In \textit{Proc.\ ICLR 2025}.

\bibitem[Macmillan-Scott \& Musolesi(2024)]{macmillan-scott2024}
Macmillan-Scott, T., \& Musolesi, M. (2024). (Ir)rationality and cognitive
biases in large language models. \textit{Royal Society Open Science}.

\bibitem[Oyserman \& Lee(2008)]{oyserman2008}
Oyserman, D., \& Lee, S.~W.~S. (2008). Does culture influence what and how we
think? Effects of priming individualism and collectivism. \textit{Psychological
Bulletin, 134}(2), 311--342.

\bibitem[Pagliaro(2026)]{pagliaro2026}
Pagliaro, A. (2026). Cognitive biases in large language models: A systematic
quantitative assessment and debiasing analysis. \textit{Electronics, 15}(11),
2428.

\bibitem[Park et al.(2023)]{park2023}
Park, J.~S., et al. (2023). Generative agents: Interactive simulacra of human
behavior. In \textit{Proc.\ UIST 2023}.

\bibitem[Rieger et al.(2015)]{rieger2015}
Rieger, M.~O., Wang, M., \& Hens, T. (2015). Risk preferences around the
world. \textit{Management Science, 61}(3), 637--648.

\bibitem[Rieger et al.(2017)]{rieger2017}
Rieger, M.~O., Wang, M., \& Hens, T. (2017). Estimating cumulative prospect
theory parameters from an international survey. \textit{Theory and Decision,
82}(4), 567--596.

\bibitem[Seo et al.(2025)]{seo2025}
Seo, M., Yuan, Y., \& Bu, Y. (2025). ValuesRAG: Enhancing cultural alignment
of large language models. \textit{arXiv:2501.01031}.

\bibitem[Suri et al.(2024)]{suri2024}
Suri, G., Slater, L.~R., Ziaee, A., \& Nguyen, J.~K. (2024). Do large
language models show decision heuristics similar to humans? \textit{Journal of
Experimental Psychology: General, 153}(4), 1066--1075.

\bibitem[Tao et al.(2024)]{tao2024}
Tao, Y., et al. (2024). Cultural bias and cultural alignment of large language
models. \textit{PNAS Nexus, 3}(9), pgae346.

\bibitem[Tversky \& Kahneman(1992)]{tversky1992}
Tversky, A., \& Kahneman, D. (1992). Advances in prospect theory: Cumulative
representation of uncertainty. \textit{Journal of Risk and Uncertainty, 5}(4),
287--297.

\bibitem[Wang et al.(2017)]{wang2017}
Wang, M., Rieger, M.~O., \& Hens, T. (2017). The impact of culture on loss
aversion. \textit{Journal of Behavioral Decision Making, 30}(2), 270--281.

\bibitem[Xie et al.(2018)]{xie2018}
Xie, Y., Hwang, S., \& Pantelous, A.~A. (2018). Loss aversion around the
world: Empirical evidence from pension funds. \textit{Journal of Banking \&
Finance, 88}, 52--62. https://doi.org/10.1016/j.jbankfin.2017.11.007
Table values used for per-country INTRA parameter transcription.

\bibitem[Yee \& Koh(2026)]{yee2026}
Yee, B., \& Koh, P. (2026). Calibrating behavioral parameters with large
language models. \textit{arXiv:2602.01022} [econ.GN]. Concurrent independent
work on profile-based calibration of behavioral parameters to generic human
benchmarks; it does not test region-specific calibration or cross-country
heterogeneity.

\end{thebibliography}

\end{document}
